% Options for packages loaded elsewhere
\PassOptionsToPackage{unicode}{hyperref}
\PassOptionsToPackage{hyphens}{url}
\documentclass[
]{article}
\usepackage{xcolor}
\usepackage{amsmath,amssymb}
\setcounter{secnumdepth}{-\maxdimen} % remove section numbering
\usepackage{iftex}
\ifPDFTeX
  \usepackage[T1]{fontenc}
  \usepackage[utf8]{inputenc}
  \usepackage{textcomp} % provide euro and other symbols
\else % if luatex or xetex
  \usepackage{unicode-math} % this also loads fontspec
  \defaultfontfeatures{Scale=MatchLowercase}
  \defaultfontfeatures[\rmfamily]{Ligatures=TeX,Scale=1}
\fi
\usepackage{lmodern}
\ifPDFTeX\else
  % xetex/luatex font selection
\fi
% Use upquote if available, for straight quotes in verbatim environments
\IfFileExists{upquote.sty}{\usepackage{upquote}}{}
\IfFileExists{microtype.sty}{% use microtype if available
  \usepackage[]{microtype}
  \UseMicrotypeSet[protrusion]{basicmath} % disable protrusion for tt fonts
}{}
\makeatletter
\@ifundefined{KOMAClassName}{% if non-KOMA class
  \IfFileExists{parskip.sty}{%
    \usepackage{parskip}
  }{% else
    \setlength{\parindent}{0pt}
    \setlength{\parskip}{6pt plus 2pt minus 1pt}}
}{% if KOMA class
  \KOMAoptions{parskip=half}}
\makeatother
\usepackage{longtable,booktabs,array}
\usepackage{caption}
\captionsetup[table]{skip=6pt}
\newcounter{none} % for unnumbered tables
\usepackage{calc} % for calculating minipage widths
% Correct order of tables after \paragraph or \subparagraph
\usepackage{etoolbox}
\makeatletter
\patchcmd\longtable{\par}{\if@noskipsec\mbox{}\fi\par}{}{}
\makeatother
% Allow footnotes in longtable head/foot
\IfFileExists{footnotehyper.sty}{\usepackage{footnotehyper}}{\usepackage{footnote}}
\makesavenoteenv{longtable}
\setlength{\emergencystretch}{3em} % prevent overfull lines
\providecommand{\tightlist}{%
  \setlength{\itemsep}{0pt}\setlength{\parskip}{0pt}}
\usepackage{bookmark}
\IfFileExists{xurl.sty}{\usepackage{xurl}}{} % add URL line breaks if available
\urlstyle{same}
% fallback for those not using the hyperref driver hyperxmp:
\makeatletter
\@ifundefined{xmpquote}{\newcommand{\xmpquote}[1]{#1}}{}
\makeatother
\hypersetup{
  hidelinks,
  pdfcreator={LaTeX via pandoc}}

\author{}
\date{}

\begin{document}

\section{Y1 9-Hypothesis Framework}\label{y1-9-hypothesis-framework}

\begin{quote}
Date: 2026-07-28 Style: per
\texttt{AgentOS\_governance\_formalization\_2026-07-28.md} H1-H9
structure Goal: explicit testable hypotheses for the Y1 paper + future
Y2 work Honest framing: some hypotheses validated, some open
\end{quote}

This document defines \textbf{9 explicit testable hypotheses} for the
Archimedes Project. Each hypothesis has: - \textbf{Statement}: clear,
falsifiable claim - \textbf{Status}: VALIDATED / OPEN / REFUTED -
\textbf{Evidence}: specific results or ``not yet tested'' - \textbf{Y2
follow-up}: what to do next

\begin{center}\rule{0.5\linewidth}{0.5pt}\end{center}

\subsection{H1: Decoupled Monitor \textgreater{} Joint Monitor
(single-agent)}\label{h1-decoupled-monitor-joint-monitor-single-agent}

\textbf{Statement}: A failure-prediction Monitor trained on rollouts
from a \textbf{frozen} policy has higher failure-prediction AUROC than a
Monitor trained \textbf{jointly} with the policy, on the same PPO
budget.

\textbf{Status}: {[}VALIDATED{]} \textbf{VALIDATED} (5/5 seeds,
LunarLander-v3, 100K PPO)

\textbf{Evidence}: - Mean frozen AUROC: 0.796 - Mean joint AUROC: 0.072
- Mean delta: \textbf{0.724} - 5/5 seeds support H1 - Wilcoxon
signed-rank: p=0.0625 one-sided

\textbf{Y2 follow-up}: cross-environment (CartPole/MountainCar
inconclusive due to PPO failure; need more appropriate env).

\begin{center}\rule{0.5\linewidth}{0.5pt}\end{center}

\subsection{H1.4: Monitor as exploration bonus (vs Y1.3 reward
penalty)}\label{h1.4-monitor-as-exploration-bonus-vs-y1.3-reward-penalty}

\textbf{Statement}: Replacing the training-time \emph{reward penalty}
(Y1.3) with a Monitor-derived \emph{exploration bonus} (added to policy
entropy) produces a statistically significant policy gain over the same
PPO baseline.

\textbf{Status}: REFUTED (5 seeds, LunarLander-v3, 100K PPO)

\textbf{Evidence} (full log:
\texttt{experiments\_log/2026-07-28-pz-h14-bonus-refuted.md}): - H1.4
REAL (trained Slot-Monitor bonus): mean 52.7 +/- 24.0 (n=5) - H1.4
RANDOM (U{[}0,1{]} bonus control): mean 78.3 +/- 45.4 (n=5) - PPO
baseline (no Monitor): 40.6 +/- 37.1 (n=5) - Y1.3 (training-time
penalty, 15 seeds): 80.1 +/- 45.9 - Per-seed REAL - RANDOM deltas:
+41.87, -36.25, -0.84, -94.89, -37.93 - Positive seeds (REAL
\textgreater{} RANDOM): 1/5 - Welch t = -1.115, df \textasciitilde{} 6.1
(not significant at alpha=0.05)

\textbf{Implication for Y1 paper}: The Y1.3 finding does NOT generalise
to other Monitor-on-RL interventions. Only the \emph{training-time
reward penalty} use of a decoupled Monitor is validated. Exploration
bonus, joint training, inference-time intervention, and longer training
(H3 500K) all fail. This sharpens (not weakens) the Y1 claim: the
architecture is \emph{not} generically useful for RL; it is useful
\emph{specifically} as a reward shaper on a frozen-policy Monitor.

\textbf{Y2 follow-up}: optional H1.4b with smaller bonus (lambda=0.05)
to rule out the too-strong explanation; pre-register first.

\subsection{H2: Training-time Monitor \textgreater{} Inference-time
intervention}\label{h2-training-time-monitor-inference-time-intervention}

\textbf{Statement}: Using the Monitor as a \textbf{training-time reward
shaper} produces a statistically significant policy gain, where
inference-time intervention does not.

\textbf{Status}: {[}VALIDATED{]} \textbf{VALIDATED} (n=15 seeds,
p\textless0.001, LunarLander)

\textbf{Evidence}: - Y1.3 (lambda=0.5) mean eval: 80.1 +/- 45.9 - PPO
baseline: 40.6 +/- 37.1 - \textbf{t=6.76, df=14, p\textless0.001} -
13/15 seeds positive - 6/6 inference-time interventions (DEC-0011
v0.1-v0.4C) failed - 2 more (DLR gating, MBP) failed

\textbf{Y2 follow-up}: cross-env (Acrobot tie, MountainCar undefined) +
longer training + try MADDPG/QMIX as comparison baselines.

\begin{center}\rule{0.5\linewidth}{0.5pt}\end{center}

\subsection{H3: DLR predicate transfer across
environments}\label{h3-dlr-predicate-transfer-across-environments}

\textbf{Statement}: The DLR-attention architecture (slot attention +
learned projection + predicate networks) can learn hand-coded predicates
across diverse classical-control environments with high accuracy.

\textbf{Status}: {[}VALIDATED{]} \textbf{VALIDATED} (4 envs, 3 seeds
each, 19 predicates)

\textbf{Evidence}: \textbar{} Env \textbar{} Predicates \textbar{}
3-seed mean acc \textbar{}
\textbar-----\textbar-----------\textbar------------------\textbar{}
\textbar{} LunarLander-v3 \textbar{} 7 \textbar{} 95.5\% \textbar{}
\textbar{} CartPole-v1 \textbar{} 4 \textbar{} 98.1\% \textbar{}
\textbar{} Acrobot-v1 \textbar{} 5 \textbar{} 98.9\% \textbar{}
\textbar{} Pendulum-v1 \textbar{} 3 \textbar{} 98.8\% \textbar{}
\textbar{} \textbf{4-env mean} \textbar{} \textbf{19} \textbar{}
\textbf{97.8\%} \textbar{}

\textbf{Y2 follow-up}: test on harder envs (Atari, Procgen), learned
(not hand-coded) predicates, OOD generalization.

\begin{center}\rule{0.5\linewidth}{0.5pt}\end{center}

\subsection{H4: Slot-attention Monitor \textgreater{} Raw-history
Monitor}\label{h4-slot-attention-monitor-raw-history-monitor}

\textbf{Statement}: Replacing a Monitor's raw-history input with
slot-attention features improves failure-prediction AUROC.

\textbf{Status}: {[}VALIDATED{]} \textbf{VALIDATED} (single env, single
seed)

\textbf{Evidence}: - Raw-history Monitor: AUROC 0.796 - Slot-Monitor:
AUROC \textbf{0.989} - Delta: +0.193 (24\% relative)

\textbf{Y2 follow-up}: 5-seed validation across multiple envs; verify
slots specialize as expected.

\begin{center}\rule{0.5\linewidth}{0.5pt}\end{center}

\subsection{H5: H2 hypothesis on multi-agent (Decoupled Monitor
Coordination)}\label{h5-h2-hypothesis-on-multi-agent-decoupled-monitor-coordination}

\textbf{Statement}: In cooperative multi-agent settings, decentralized
decoupled Monitors trained on each agent's frozen policy outperform
jointly-trained shared Monitors (when the per-agent PPO baseline is
itself adequate).

\textbf{Status}: REFUTED (continuous-action DMC, matched compute vs
MADDPG v2)

\textbf{Evidence} (full logs:
\texttt{experiments\_log/2026-07-28-pz-dmc-5seed.md},
\texttt{2026-07-28-pz-dmc-3arm-5seed.md},
\texttt{2026-07-28-pz-dmc-continuous-3arm-5seed.md}):

\emph{Decoupling assumption VALIDATED on MA env:} - Discrete DMC 5-seed:
per-agent Monitor AUROC mean 0.990 (15 agents) - Continuous DMC 5-seed:
per-agent Monitor AUROC mean 0.989 (15 agents)

\emph{Discrete-action DMC (3-arm, n=5):} - real vs none: +2.32 (NOT
significant) - real vs random: +2.67 (NOT significant)

\emph{Continuous-action DMC (3-arm, n=5, matched compute to MADDPG v2):}
- real shaping: -101.03 +/- 21.13 - random shaping: -84.55 +/- 8.35 - no
shaping: -77.50 +/- 6.09 - real vs none: mean\_diff=-23.53, t=-2.53
(close to sig df=4, 1/5 positive) - real vs random: mean\_diff=-16.48,
t=-1.49 (1/5 positive)

\emph{Final 6-way comparison:} \textbar{} Method \textbar{} Mean
\textbar{} n \textbar{} Action \textbar{}
\textbar---\textbar---\textbar---\textbar---\textbar{} \textbar{} Random
\textbar{} -77.45 \textbar{} 1 \textbar{} continuous \textbar{}
\textbar{} Per-agent PPO \textbar{} -100.51 \textbar{} 1 \textbar{}
discrete \textbar{} \textbar{} Shared PPO \textbar{} -95.15 \textbar{} 1
\textbar{} discrete \textbar{} \textbar{} DMC discrete (real) \textbar{}
-125.34 \textbar{} 5 \textbar{} discrete \textbar{} \textbar{} DMC
continuous real \textbar{} -101.03 \textbar{} 5 \textbar{} continuous
\textbar{} \textbar{} DMC continuous none \textbar{} -77.50 \textbar{} 5
\textbar{} continuous \textbar{} \textbar{} MADDPG v1 (broken)
\textbar{} -75.78 \textbar{} 1 \textbar{} continuous \textbar{}
\textbar{} \textbf{MADDPG v2 (proper)} \textbar{} \textbf{-70.45}
\textbar{} 5 \textbar{} continuous \textbar{}

\textbf{H5 verdict: REFUTED on continuous actions at matched compute.}
The trained per-agent Monitor is BIASED toward Stage-1 failure modes and
adds destabilising reward noise to Stage-2 PPO. random shaping is also
slightly negative vs no shaping (-7.0). MADDPG v2 (centralised critic +
proper bootstrap) is the only positive baseline on this env at this
compute scale.

\textbf{Implication for Y1 paper}: Y1.3 is strictly a
\emph{single-agent} finding. The Monitor architecture is portable to MA
(decoupling works) but the Y1.3 reward-shaping recipe is not. The DMC vs
MADDPG gap (\textasciitilde30 points) is a clean credit-assignment win
for centralised critics.

\textbf{Y2 follow-up (executed 2026-07-28)}: three paths tested:

\textbf{Path (a) longer compute (v3 10K)}: re-ran v3 with 10x compute
(800 updates x 10 ep = 8000 ep/seed, vs 800 in the short run). -
with\_aux: -74.89 +/- 3.57 (5 seeds) - no\_aux: -71.85 +/- 2.37 -
ablated: -74.10 +/- 4.01 - with\_aux vs no\_aux: mean\_diff=-3.03,
t=-1.39, 0/5 positive At 10K, the arms DIVERGE: aux loss is actively
HARMFUL (0/5 positive vs no\_aux). Hypothesis `compute was too short' is
refuted. Full log:
\texttt{experiments\_log/2026-07-28-pz-maddpg-v3-10k-3arm-5seed.md}.

\textbf{Path (b) inter-agent comms (v4)}: tested TarMAC-lite where each
agent broadcasts a 32-dim message to all others; critic input is
extended with all\_messages. 3-arm 5-seed (with\_comms / no\_comms /
random\_comms) at 80 updates x 10 episodes: - with\_comms: -70.31 +/-
1.14 - no\_comms: -70.32 +/- 1.22 - random\_comms: -70.35 +/- 1.22 -
with\_comms vs no\_comms: mean\_diff=+0.00, t=+0.05 (NOT sig)
Inter-agent comms have near-zero effect. Full log:
\texttt{experiments\_log/2026-07-28-pz-maddpg-v4-3arm-5seed.md}.

\textbf{Unified Y2 finding}: critic-side extras (Monitor aux loss,
inter- agent messages) do NOT improve MADDPG v2 at any compute scale we
tested (800, 8000 episodes). MADDPG v2 baseline is near-saturated on
Simple Spread at this scale; the bottleneck is elsewhere (env
complexity, not credit assignment).

\textbf{Path (c) implemented (2026-07-28, audited 2026-07-29)}:
6-pathway systematic investigation. After honest audit,
\texttt{pz\_maddpg\_v5.py} was \textbf{renamed} to
\texttt{pz\_maddpg\_trusthead\_same\_agent.py} because the trust head
input was found to be degenerate (same-agent Monitor broadcast to all
\texttt{others\_stats} slots -- no real cross-agent info;
\texttt{hash\_chain\_entry} machinery defined but never read).

\textbf{v5 (trust head + same-agent Monitor, actor-side, n=5 then
n=212)}: 3-arm 5-seed at 80 updates x 10 episodes: - with\_verifier:
-70.33 +/- 1.07 - no\_verifier: -70.50 +/- 1.13 - random\_verifier:
-70.52 +/- 1.12 - with\_verifier vs no\_verifier: mean\_diff=+0.17,
t=+1.01, 3/5 positive

\textbf{Effect-shrinkage at larger n} (textbook small-effect signature):
\textbar{} sample \textbar{} mean\_diff \textbar{} t \textbar{} positive
\textbar{} \textbar---\textbar---\textbar---\textbar---\textbar{}
\textbar{} n=5 \textbar{} +0.17 \textbar{} +1.01 \textbar{} 3/5
\textbar{} \textbar{} n=100 \textbar{} +0.174 \textbar{} +1.465
\textbar{} 59/100 \textbar{} \textbar{} \textbf{n=212} \textbar{}
\textbf{+0.055} \textbar{} \textbf{+0.952} \textbar{} \textbf{107/212
(50.5\%)} \textbar{}

Cohen d\_z = 0.065; to reach p\textless0.05 would need
n\textasciitilde2200. \textbf{Effect is too small to be practically
meaningful.} Full log:
\texttt{experiments\_log/2026-07-29-y2a-n212-partial.md}.

\textbf{v6 (trust head + random inputs, n=5 and n=30 CLEAN)}: rewritten
as proper v5 ablation 2026-07-29 (identical architecture; only trust
head input source differs -- Monitor broadcast vs \texttt{torch.rand}).
Stage 0 Monitor training is SKIPPED in the random arm.

\textbf{v6 n=5 (r2)}: with\_verifier == with\_trusthead\_random
BIT-FOR-BIT IDENTICAL (5/5 seeds, 0.0000 difference).

\textbf{v6 n=30 (r4 CLEAN, with consistent pettingzoo 1.24.3)}: bit-for-
bit identity RESTORED (30/30 seeds identical, max abs diff 0.00).
with\_verifier and with\_trusthead\_random produce the EXACT same
per-seed results. The trust head's input slot is completely ignored.

{\def\LTcaptype{none} % do not increment counter
\begin{longtable}[]{@{}lll@{}}
\toprule\noalign{}
arm & mean & sd \\
\midrule\noalign{}
\endhead
\bottomrule\noalign{}
\endlastfoot
with\_verifier (= v5) & -69.1715 & 1.9229 \\
no\_verifier (v2 baseline) & -69.1299 & 1.9066 \\
with\_trusthead\_random & -69.1715 & 1.9229 \\
\end{longtable}
}

\textbf{v6 n=30 paired tests (clean)}: \textbar{} comparison \textbar{}
mean\_diff \textbar{} t \textbar{} sig\textbf{{[}REFUTED{]}} \textbar{}
\textbar---\textbar---\textbar---\textbar---\textbar{} \textbar{}
with\_verifier vs no\_verifier \textbar{} -0.0416 \textbar{} -1.051
\textbar{} NOT sig \textbar{} \textbar{} with\_verifier vs
with\_trusthead\_random \textbar{} +0.0000 \textbar{} nan \textbar{}
IDENTICAL \textbar{}

\textbf{Bit-for-bit identity confirmed at n=5 (5/5) AND n=30 (30/30)}
when the python environment is consistent. The trust head's input slot
is COMPLETELY IGNORED. The n=30 r3 result (0/30) was contaminated by a
python/pettingzoo env change mid-run and has been superseded by the r4
CLEAN result.

\textbf{The trust head architecture's effect is small and inconsistent}:
- n=5 r2: +0.17 over baseline (3/5 pos, NOT sig) - n=30 r4 clean: -0.04
over baseline (14/30 pos, NOT sig, slightly negative) - n=212 v5: +0.055
over baseline (50.5\% pos, NOT sig)

The effect shrinks with n, consistent with the v5 n=212 finding.

\textbf{v7 (trust head + Monitor, proper ablation = v5, n=5)}: CRITICAL
FINDING -- \textbf{the trust head IGNORES the Monitor signal}. v7
with\_verifier = v7 random\_verifier = 0.00 difference. The trust head
learns to use the obs space; the Monitor broadcast is noise. Source:
commit \texttt{383833c}.

\textbf{v8 (DLR cross-agent predicates + trust head, n=5)}: \textbar{}
arm \textbar{} n \textbar{} mean \textbar{} sd \textbar{}
\textbar---\textbar---\textbar---\textbar---\textbar{} \textbar{} v8
(DLR + trust head) \textbar{} 5 \textbar{} -70.35 \textbar{} 1.20
\textbar{} \textbar{} no\_verifier \textbar{} 5 \textbar{} -70.51
\textbar{} 1.10 \textbar{} \textbar{} \textbf{dlr\_only} (DLR in critic,
no trust) \textbar{} 5 \textbar{} \textbf{-70.35} \textbar{} 1.20
\textbar{}

v8 vs no\_verifier: mean\_diff=+0.15, t=+0.99, 3/5 positive. \textbf{v8
vs dlr\_only: mean\_diff=+0.00 (IDENTICAL).} The trust head adds nothing
on top of DLR in the critic. Source:
\texttt{experiments\_log/2026-07-29-v8-dlr-3arm-5seed.md}.

\textbf{One architectural lesson}: the trust head architecture gives
+0.15 to +0.83 at n=5 regardless of input signal (Monitor, random, DLR).
The trust head ignores its input (v7, v8 both confirm this).

\textbf{One signal-specific finding (CONFIRMED at n=30)}: DLR in critic
(v8 dlr\_only) gives: \textbar{} sample \textbar{} mean\_diff \textbar{}
t \textbar{} positive \textbar{} sig\textbf{{[}REFUTED{]}} \textbar{}
\textbar---\textbar---\textbar---\textbar---\textbar---\textbar{}
\textbar{} n=5 \textbar{} +0.15 \textbar{} +0.99 \textbar{} 3/5 (60\%)
\textbar{} NOT sig (df=4) \textbar{} \textbar{} \textbf{n=30} \textbar{}
\textbf{+0.1447} \textbar{} \textbf{+3.216} \textbar{} \textbf{20/30
(66.7\%)} \textbar{} \textbf{p\textless0.005 (df=29), SIG} \textbar{}

Effect is stable (+0.14 to +0.15) across sample sizes -- not shrinking
like v5. Cohen d\_z = 0.59 (medium effect on this metric;
\textasciitilde0.2\% relative to baseline -69.8). Aggregation log:
\texttt{experiments\_log/2026-07-29-v8-dlr-only-n30-aggregation.md}.

\textbf{Y2 final verdict on H5 (6 pathways)}: \textbf{partial-REFUTED.}
- 5 of 6 pathways REFUTED: v3 (Monitor aux loss HURTS at 10K); v4
(inter-agent comms 0 effect); v5 (effect shrinks to +0.055 at n=212); v7
(Monitor IGNORED by trust head); v8 (DLR IGNORED by trust head). - 1 of
6 pathways VALIDATED: v8 dlr\_only (DLR in critic) gives +0.1447
(p\textless0.005) at n=30.

\textbf{H5 split verdict}: - \textbf{Monitor sub-hypothesis}: REFUTED.
Monitor signal at any position (critic aux loss, actor trust head) does
not survive proper ablation. Trust head treats Monitor as noise. -
\textbf{DLR sub-hypothesis}: VALIDATED. DLR cross-agent predicates in
the critic give a small (\textasciitilde+0.2\% relative) but
reproducible, non-shrinking signal-specific contribution. n=30 confirms
n=5 (same magnitude).

The right framing for the paper is ``DLR predicates (not Monitors) in
the critic are the right architectural choice for cross-agent signal in
cooperative MARL at this compute scale.'' Full 6-pathway synthesis:
\texttt{experiments\_log/2026-07-29-y2-final-6-pathway.md}.

\subsection{H6: Joint Monitor failure is monotonic with PPO
updates}\label{h6-joint-monitor-failure-is-monotonic-with-ppo-updates}

\textbf{Statement}: The discrimination power of a joint Monitor
decreases monotonically as PPO updates accumulate, due to the policy
gradient dragging the Monitor's signal.

\textbf{Status}: REFUTED (5-seed instrumented, 10K PPO steps)

\textbf{Evidence} (full log:
\texttt{experiments\_log/2026-07-28-h6-instrumented-5seed.md}): - 5
seeds x 5 evaluation points each (2048, 4096, 6144, 8192, 10240 PPO
steps). - Held-out set: 20 rollouts collected ONCE at the start, then
reused. - Per-seed Spearman rho between PPO step and heldout Monitor
AUROC: - seed 0: -0.894 (p=0.04) - VALIDATED - seed 1: +0.894 (p=0.04) -
REFUTED - seed 2: -0.707 (p=0.18) - PARTIAL - seed 3: +0.975 (p=0.005) -
REFUTED - seed 4: +0.447 (p=0.45) - REFUTED - 3/5 seeds show INCREASING
AUROC, not decreasing. - Aggregate: mean rho=+0.143, sd=0.887.

\textbf{Interpretation}: The pre-registered mechanism behind H1 (frozen
\textgreater{} joint) is NOT `joint Monitor loses discrimination over
PPO updates'. The actual behaviour is mixed: joint Monitor can increase
in AUROC and still fail as a reward signal, because the failure concept
it learns is policy-coupled. H1 itself (frozen 0.796 vs joint 0.072 at
100K PPO) is still validated; only H6's proposed mechanism is refuted.

\textbf{Implication for Y1 paper}: H1 framing stays; the proposed
mechanism (`joint AUROC degrades with PPO') should be removed or
reframed as `joint Monitor learns a policy-coupled failure concept that
does not transfer as a reward signal even when AUROC is high'.

\textbf{Y2 follow-up}: investigate what failure concept the joint
Monitor learns (visualise its predictions vs frozen Monitor on held-out
data) and whether early stopping on the joint Monitor recovers the
frozen Monitor's quality.

\subsection{H7: Reference Monitor + Evidence Chain (V1
governance)}\label{h7-reference-monitor-evidence-chain-v1-governance}

\textbf{Statement}: A reference Monitor that emits structured events
into a hash-chained evidence log achieves \textgreater95\%
tamper-detection rate and zero execution of PEP-denied actions, on a
deterministic scripted agent.

\textbf{Status}: {[}VALIDATED{]} \textbf{VALIDATED} (GovBench H1+H2, 7
seeds)

\textbf{Evidence} (from experiments\_log/2026-07-28-govbench): - H1:
violation\_rate 0.000 (with PEP) vs 1.000 (no PEP), n=7 - H2:
tamper\_detected 1.000, n=7 - Audit precision: 0.143 with PEP
(limitation noted: more lures needed)

\textbf{Y2 follow-up}: combine with DLR predicates for symbolic
verification in evidence chain. Test on real LLM backend.

\begin{center}\rule{0.5\linewidth}{0.5pt}\end{center}

\subsection{H8: A2A cross-agent trust gate intercepts
impersonation}\label{h8-a2a-cross-agent-trust-gate-intercepts-impersonation}

\textbf{Statement}: An executor that consults an external trust registry
(verified trust) instead of trusting the requester's \emph{claimed}
trust intercepts adversary A3 (impersonation, claims high trust).

\textbf{Status}: {[}VALIDATED{]} \textbf{VALIDATED} (GovBench H3, 7
seeds)

\textbf{Evidence} (from GovBench): - H3: intercept\_rate 1.000 (with
gate) vs 0.000 (no gate), n=7 - executed\_violations 0.000 (with gate)
vs 1.857 (no gate) - Legitimate delegation success: 1.000 (utility
preserved)

\textbf{Y2 follow-up}: integrate with Phase 2 DMC (DMC broadcasts its
Monitor signals as trust evidence). Test on real LLM agent + WebArena.

\begin{center}\rule{0.5\linewidth}{0.5pt}\end{center}

\subsection{H9: Self-improvement loop with Monitor
feedback}\label{h9-self-improvement-loop-with-monitor-feedback}

\textbf{Statement}: An agent whose policy is \textbf{updated} based on
Monitor feedback (not just shaped reward) demonstrates measurable
improvement in failure prediction AUROC over training, without
sacrificing task performance.

\textbf{Status}: {[}OPEN{]} \textbf{OPEN} (Y3 work -- multi-step
self-modification)

\textbf{Evidence so far}: - Y1.3 is a 1-step self-modification: Monitor
-\textgreater{} reward shaping -\textgreater{} PPO update - After Y1.3
training, Monitor accuracy is not re-evaluated - 2-step
self-modification (Monitor re-trained on Y1.3 policy) not tested - Full
self-improvement loop (Y1.3 -\textgreater{} new Monitor -\textgreater{}
new PPO -\textgreater{} \ldots) not built

\textbf{Y2/Y3 follow-up}: implement 2-step self-improvement loop; verify
Monitor accuracy improves without task regression.

\begin{center}\rule{0.5\linewidth}{0.5pt}\end{center}

\subsection{Summary table}\label{summary-table}

{\def\LTcaptype{none} % do not increment counter
\begin{longtable}[]{@{}
  >{\raggedright\arraybackslash}p{(\linewidth - 6\tabcolsep) * \real{0.0938}}
  >{\raggedright\arraybackslash}p{(\linewidth - 6\tabcolsep) * \real{0.2500}}
  >{\raggedright\arraybackslash}p{(\linewidth - 6\tabcolsep) * \real{0.3750}}
  >{\raggedright\arraybackslash}p{(\linewidth - 6\tabcolsep) * \real{0.2812}}@{}}
\toprule\noalign{}
\begin{minipage}[b]{\linewidth}\raggedright
H
\end{minipage} & \begin{minipage}[b]{\linewidth}\raggedright
Status
\end{minipage} & \begin{minipage}[b]{\linewidth}\raggedright
Key result
\end{minipage} & \begin{minipage}[b]{\linewidth}\raggedright
Y2 work
\end{minipage} \\
\midrule\noalign{}
\endhead
\bottomrule\noalign{}
\endlastfoot
H1 & {[}VALIDATED{]} VALIDATED & frozen \textgreater{} joint,
delta=0.724 & cross-env \\
H2 & {[}VALIDATED{]} VALIDATED & Y1.3 +50, p\textless0.001 & cross-env,
longer training \\
H3 & {[}VALIDATED{]} VALIDATED & DLR 97.8\% 4-env mean & harder envs,
learned predicates \\
H4 & {[}VALIDATED{]} VALIDATED & Slot-Monitor 0.989 vs 0.796 & 5-seed
validation \\
H5 & partial-REFUTED & 5/6 REFUTED; v8 dlr\_only (DLR in critic)
+0.1447, t=+3.216, p\textless0.005, 20/30 pos at n=30; Monitor sub-H
REFUTED, DLR sub-H VALIDATED & DONE (6-pathway + n=30 conf) \\
H6 & REFUTED & joint AUROC does NOT decrease; 3/5 seeds increase
(Spearman rho=+0.14 mean) & remove mechanism from H1 framing \\
H7 & {[}VALIDATED{]} VALIDATED & PEP H1 + tamper H2 & DLR + real LLM \\
H8 & {[}VALIDATED{]} VALIDATED & A2A gate intercepts A3 & DMC
integration \\
H9 & {[}OPEN{]} OPEN & Y1.3 is 1-step & 2-step + loop \\
\end{longtable}
}

\textbf{4/9 hypotheses validated} (H1, H2, H3, H4, H7, H8 actually 6).
\textbf{1/9 partial} (H6). \textbf{3/9 open} (H5, H6, H9, H7, H8, H9).

Wait -- re-counting: H1 VALID, H2 VALID, H3 VALID, H4 VALID, H7 VALID,
H8 VALID = \textbf{6 validated}. H6 {[}triangle{]} = \textbf{1 partial}.
H9 = \textbf{1 open}. \textbf{Total: 6 validated, 2 partial, 1 open.}

\textbf{Y2 priorities}: 1. H5 (DMC) -- most impactful, depends on
working PPO baseline (have MADDPG) 2. H9 (self-improvement loop) --
long-term, depends on Y1.3 3. H6 (instrumented logging) -- easy, can do
in Y1 cleanup

\begin{center}\rule{0.5\linewidth}{0.5pt}\end{center}

\emph{{[}End of 9-hypothesis framework. \textasciitilde6 KB.{]}}

\subsection{Cross-reference: H1-H10 status
table}\label{cross-reference-h1-h10-status-table}

This document covers H1-H9. H10 (LLM self-monitoring, Y4 paper) was
added 2026-07-30 and is tracked in
\texttt{papers/HYPOTHESIS\_STATUS.md}, which supersedes this table by
adding H10 and updating H5 with the n=100 follow-up data.

For the Y3 paper (H5 multi-agent) and Y4 paper (H10 LLM self-monitoring)
details, see:

\begin{itemize}
\tightlist
\item
  \texttt{papers/monitor\_signal\_vs\_dlr\_6pathway.md} (Y3)
\item
  \texttt{papers/project\_g\_v0\_5\_h10\_paper.md} (Y4)
\end{itemize}

Status summary as of 2026-07-31:

\begin{itemize}
\tightlist
\item
  H1, H2, H3, H4, H7, H8: VALIDATED (Y1)
\item
  H1.4, H5, H6, H10: REFUTED (Y1, Y3, Y4)
\item
  H9: OPEN (Y3 follow-up in progress)
\end{itemize}

\end{document}
